% ─── Chapter 3: Challenge Model & Data Provenance ──────────────────
\chapter{Challenge Model \& Data Provenance}
\label{ch:challenge}

\section{Challenge Definition}

Instead of playing through the entire Pok\'{e}mon game world, the project
focuses on \textbf{battle sequences}.  The challenge is modelled as a graph of
battles: each node represents a battle against a specific trainer team, and
edges represent progression.

In the simplest form---currently implemented---the challenge is a linear list
of opponents.  The agent must beat all trainers in order, without losing.  If
the agent loses a battle, the run ends and restarts from the beginning.

\subsection{Resource Management Between Battles}

After each battle, the agent's resources are \emph{partially} restored:

\begin{itemize}
  \item HP and status conditions are fully healed.
  \item Power Points (PP, move usage limits) are restored.
  \item \textbf{Fainted Pok\'{e}mon remain lost} (Nuzlocke rule).
\end{itemize}

Usually, the agent may build a new team before each fight.  There are
exceptions where the agent must battle through a sequence of opponents (up to
five) before being allowed to change teams again.  The amount of information
available to the team builder---complete, partial, or none---serves as a knob
for testing agent competency.

\section{Milestone Roadmap}

Due to the large scope of the project, four progressive milestones have been
defined.  For each milestone, a model is trained from scratch and evaluated
against the gauntlet to validate competency.

\begin{center}
\small
\begin{tabularx}{\textwidth}{c l X}
  \toprule
  \textbf{\#} & \textbf{Milestone} & \textbf{Goal} \\
  \midrule
  1 & Proof of Concept &
    The agent uses a \emph{fixed team} (6 Pok\'{e}mon, reused every battle)
    against \emph{random opponents}.  Tests whether the model can learn to
    use a known team against unseen opponents. \\
  2 & Random Battles &
    Both the agent and opponent use \emph{randomly generated teams}.  Tests
    competency under high entropy, where each battle is novel. \\
  3 & Team Building &
    Model~1 learns to compose optimal teams given opponent information;
    Model~2 learns to battle with those teams.  The team builder must
    construct a team that can beat at least five different opponents. \\
  4 & Gauntlet Master &
    Full gauntlet completion under Nuzlocke constraints.  The ultimate
    objective. \\
  \bottomrule
\end{tabularx}
\end{center}

\textbf{Current status:} Milestone~1 is in progress.  The training
infrastructure is fully functional and multiple architectural improvements have
been validated (see Chapter~\ref{ch:experiments}), but gauntlet-level
competency has not yet been demonstrated.

\section{Data Provenance}

All data used in this project is derived from publicly available sources or
generated programmatically.  No personal data is involved.

\subsection{BDSP Trainer Dataset}

\begin{center}
\small
\begin{tabularx}{\textwidth}{l l X}
  \toprule
  \textbf{File} & \textbf{Format} & \textbf{Content} \\
  \midrule
  \texttt{data/BDSP Trainer Data - v1.1.csv} & CSV &
    Complete roster of 297 BDSP trainers with Pok\'{e}mon species, levels,
    movesets, held items, and abilities. \\
  \texttt{data/bdsp\_gauntlet\_order.json} & JSON &
    Ordered list of 706 trainers defining the gauntlet sequence, sorted by
    ascending difficulty (level, team size).  Includes metadata about
    breakpoints and battle counts. \\
  \texttt{data/bdsp\_trainers.json} & JSON &
    Trainer team specifications in Showdown-compatible format.  Used to
    instantiate specific opponents during training and evaluation. \\
  \bottomrule
\end{tabularx}
\end{center}

\subsection{Vocabulary Files}

Categorical entities (species, items, abilities) are mapped to integer IDs via
stable dictionary lookups stored in JSON:

\begin{center}
\small
\begin{tabularx}{0.8\textwidth}{l r l}
  \toprule
  \textbf{Entity} & \textbf{Vocabulary Size} & \textbf{Source} \\
  \midrule
  Species  & 1\,151 & All Gen\,8 species \\
  Items    & 1\,110 & All available held items \\
  Abilities & 216   & All available abilities \\
  \bottomrule
\end{tabularx}
\end{center}

These vocabularies were generated programmatically from Pok\'{e}mon Showdown's
data files and validated for complete coverage over the BDSP trainer dataset
and validation manifests (zero unknown entities).

\subsection{Validation Team Manifests}

For reproducible evaluation, a set of frozen team configurations is maintained:

\begin{itemize}
  \item \texttt{data/validation/gen8\_random\_battle\_team\_pairs.json}:
        20 teams organised as 10 swapped pairs, yielding 40 battles per
        evaluation run.
  \item \texttt{data/validation/gen8\_random\_battle\_mirror\_teams.json}:
        Mirror-matchup teams for symmetric evaluation.
\end{itemize}

These teams are evaluated against both \texttt{RandomPlayer} and
\texttt{SimpleHeuristicsPlayer} opponents, providing a fixed benchmark for
comparing checkpoints.

\section{Data Preprocessing Pipeline}

Raw data undergoes the following transformations before reaching the agent:

\begin{enumerate}
  \item \textbf{Trainer data loading:} CSV is parsed into an internal
        representation; team JSON files are loaded and converted to
        Showdown-format team strings.
  \item \textbf{Team generation:} For random battles, teams are sampled from
        the Gen\,8 random battle pool with seed control for reproducibility.
  \item \textbf{Observation construction:} Each turn, the raw battle state
        from \texttt{poke-env} is encoded into a fixed-size tensor
        ($13 \times 164$ floats) plus categorical ID vectors for species,
        items, and abilities.  The full encoding pipeline is detailed in
        Chapter~\ref{ch:architecture} and Appendix~\ref{ch:appendix_obs}.
  \item \textbf{Action mapping:} The agent's compressed 14-action output is
        mapped to native \texttt{poke-env} \texttt{BattleOrder} objects at
        runtime (see Section~\ref{sec:action_space}).
\end{enumerate}
